% SRC: report.md §3, §4.3, §5.1-5.3 (all numbers from the recovered draft)
\section{Scene reconstruction results}

UoS Vitrine runs as independently repeatable stages (ingest, SfM, training, evaluation,
export, mesh derivation, packaging, verification), each writing a compact JSON report so
expensive stages can be inspected or repeated in isolation. The same algorithm serves every
hardware tier; profiles change only measured quantities (resolution, crop size, Gaussian
cap, iterations). Mixed stills and video are handled with separate camera models (Figure~\ref{fig:coverage}): applying
one set of intrinsics across both groups can silently warp the reconstruction. Training uses
gsplat's MCMC densification under a hard population cap, random high-resolution crops to keep
full-resolution texture in reach without rendering whole 4K frames per step, and an objective
of $0.8\,L_1 + 0.2\,(1-\mathrm{SSIM})$ because pixel-wise losses stay tolerant of blur ---
during development a blurred image held \SI{10.8}{\decibel} PSNR while SSIM collapsed to 0.013.

\begin{figure}[t]
\centering
% SRC: uos-vitrine/docs/images/nested-cinema-colmap-camera-coverage.png (Glenn Watts, UoS Vitrine repo)
\includegraphics[width=0.96\columnwidth]{nested-cinema-colmap-camera-coverage.png}
\caption{COLMAP camera coverage over the room-set: the 222 registered stills and video frames
solved as two camera models. Source: UoS Vitrine repository (G. Watts).}
\label{fig:coverage}
\end{figure}


\subsection{Laptop and workstation measurements}
On the target laptop (RTX 3060 Laptop, \SI{6}{\giga\byte}) the validated run staged 222 images,
registered all of them, and recovered two camera models and 111{,}203 sparse points. Training
completed 15{,}000 iterations in \SI{60.4}{\minute}, ending at one million Gaussians with
full degree-three spherical harmonics, \SI{1.94}{\giga\byte} peak VRAM, and held-out metrics
of \SI{25.72}{\decibel} PSNR / 0.8167 SSIM. VRAM was never the binding constraint; elapsed
time was. The RTX 5090 control reused the proven recipe (768-pixel crop, 2048-pixel source,
one-million cap, 15k iterations) and reached \SI{25.06}{\decibel} / 0.8211 SSIM in
\SI{2.4}{\minute} at \SI{1.75}{\giga\byte}, matching laptop quality about 25 times faster.

Two measurements explain why time, not memory, binds the laptop. Rasterisation cost tracks
the number of Gaussians projected into the camera frustum, not the total population
(Table~\ref{tab:frustum}): an interior camera typically observes about a third of the room,
so the middle row is the practical basis for laptop estimates. The RTX 5090 reduces a
comparable step to 12--14\,ms, at which point kernel-launch overhead flattens the
frustum--time relationship. The same six profiles span both tiers by changing only measured
quantities (Table~\ref{tab:profiles}); the workstation-archive row records the investigated
configuration, not an approved preset.

\begin{table}[t]
\centering\small
\caption{Laptop frustum occupancy dominates iteration time (1.5M Gaussians, $768^2$, SH~3).
% SRC: report.md §4.1 (lines 193-199)
}
\label{tab:frustum}
\setlength{\tabcolsep}{5pt}
\begin{tabular}{@{}rr@{}}
\toprule
Gaussians in view & Time / iteration \\
\midrule
94\% & \SI{500}{\milli\second} \\
39\% & \SI{222}{\milli\second} \\
15\% & \SI{95}{\milli\second} \\
\bottomrule
\end{tabular}
\end{table}

\begin{table}[t]
\centering\small
\caption{Profiles. Same algorithm; only measured quantities change.
% SRC: report.md §4.3 (lines 209-214)
}
\label{tab:profiles}
\setlength{\tabcolsep}{3pt}
\begin{tabular}{@{}lrrrr@{}}
\toprule
Profile & Src & Crop & Cap & Iters \\
\midrule
laptop-draft       & 1600 & 512  & 400k & 7k \\
laptop-standard    & 2048 & 768  & 1.0M & 15k \\
laptop-archive     & 2560 & 768  & 1.5M & 30k \\
workstation-draft  & 2048 & 800  & 1.0M & 7k \\
workstation-standard & 3200 & 1280 & 3.0M & 20k \\
workstation-archive & 4096 & 1600 & 6.0M & 30k \\
\bottomrule
\end{tabular}
\end{table}

\subsection{Archive-profile opacity collapse}
Our strongest tested negative result isolates crop/source scale as the decisive variable
(Table~\ref{tab:collapse}). Both healthy runs use 768/2048; both collapsed runs use 1600/4096.
Reducing the cap from six million to one million did not restore the live population, and
correcting the position-learning-rate horizon did not either. In the large-crop runs the
fraction of Gaussians above the live-opacity threshold falls catastrophically: one random
1600-pixel crop from a 4096-pixel source does not revisit each Gaussian often enough for
reconstruction gradients to counter global opacity and scale regularisation, so MCMC relocates
dead splats from a shrinking live pool. The empirical conclusion that 1600/4096 is unsafe in the
current trainer is definitive; the mechanism remains a testable hypothesis. Live exported
population therefore belongs beside PSNR and SSIM --- a nominal six-million cap says little if
98.7\% of the population is effectively transparent.

\begin{table}[t]
\centering\small
\caption{RTX 5090 controlled matrix. Crop/source scale, not cap, tracks the collapse.
% SRC: report.md §5.3 table (lines 255-261)
}
\label{tab:collapse}
\setlength{\tabcolsep}{3pt}
\begin{tabular}{@{}lrrrr@{}}
\toprule
Run & Crop/src & Cap & Live g. & PSNR \\
\midrule
3060 baseline    & 768/2048  & 1.0M & 555k (55\%)  & 25.72 \\
5090 control     & 768/2048  & 1.0M & 528k (53\%)  & 25.06 \\
5090 candidate Q1& 1024/3072 & 1.5M & 940k (63\%)  & 24.71 \\
5090 safe-cap    & 1600/4096 & 1.0M & 2.2k (0.2\%) & 17.49 \\
5090 six-million & 1600/4096 & 6.0M & 80k (1.3\%)  & 14.31 \\
\bottomrule
\end{tabular}
\end{table}

A matched-view diagnostic on held-out still \texttt{IMG\_6351} (Figure~\ref{fig:photovsrecon}) shows both splat renders soften
sofa fabric, newspaper print and equipment detail that stay sharp in the source. The detail
exists in the stills but does not survive reconstruction; the indicated next test is a
stills-only baseline before reintroducing only the video frames that close a genuine coverage
gap. The 5090 control is the honest working workstation result: it reproduces laptop quality
about 25 times faster but does not yet exceed it.

\begin{figure*}[t]
\centering
% SRC: uos-vitrine/docs/images/nested-cinema-photo-vs-reconstruction.png (Glenn Watts, UoS Vitrine repo)
\includegraphics[width=0.92\textwidth]{nested-cinema-photo-vs-reconstruction.png}
\caption{Matched-view comparison on the newspaper partition: source photograph (left) against the
splat reconstruction (right) from the same COLMAP camera. The printed text and clip edges stay
sharp in the source and soften in the render --- the detail is present in the stills but does not
survive reconstruction. Source: UoS Vitrine repository (G. Watts).}
\label{fig:photovsrecon}
\end{figure*}


\subsection{Energy and cost}
The Windows runs sample GPU power through \texttt{nvidia-smi} and integrate it over training,
costed at the recorded \pounds0.2703/kWh tariff (Table~\ref{tab:energy}). The control trains for
about a third of a penny; the collapsed safe-cap run costs nearly four times as much to produce a
worse result. These figures cover sampled GPU power, not whole-system wall energy, so they are
valid for relative trainer comparison on this workstation but are not the full institutional cost.
A wall-socket meter is the preferred boundary for the planned sustainability integration.

\begin{table}[h]
\centering\small
\caption{Sampled-GPU energy and cost per 5090 run.
% SRC: report.md §5.3 energy table (lines 277-281)
}
\label{tab:energy}
\setlength{\tabcolsep}{4pt}
\begin{tabular}{@{}lrrl@{}}
\toprule
Run & Time & Energy & Cost \\
\midrule
Control 768/2048   & \SI{2.4}{\minute} & \SI{0.013}{kWh} & \pounds0.0035 \\
Q1 1024/3072       & \SI{3.0}{\minute} & \SI{0.018}{kWh} & \pounds0.0049 \\
Safe-cap 1600/4096 & \SI{7.6}{\minute} & \SI{0.049}{kWh} & \pounds0.0132 \\
\bottomrule
\end{tabular}
\end{table}
